\documentclass[UTF8]{beamer}
\usepackage{graphicx, color}
\usepackage{algorithm2e}
\usepackage{zhspacing}
\usepackage{amsmath}

\usepackage{underscore}
\usetheme{JuanLesPins}
\usepackage{fontspec}
\setsansfont{Microsoft YaHei}

\usepackage{enumerate}

\AtBeginSection[]{
  \frame{
    \frametitle{Next}
    \tableofcontents[currentsection, subsectionstyle=show/shaded/hide]
  }
}

\AtBeginSubsection[]{
  \frame{
    \frametitle{Next}
    \tableofcontents[currentsubsection]
  }
}

\title{Perl in Bioinformatics}

\author{Gang Chen\\ chengang@bgitechsolutions.com}

\logo{\includegraphics[width=1.3cm]{bgi-logo.png}\includegraphics[width=2.5cm]{cuhklogo.png}}
\date{\today}




\begin{document}

\begin{frame}
  \frametitle{Science: An Explosion Of Bioinformatics Careers}
\begin{block}{Programming}
  Additional critical skills are required for big data careers in industry, such
   as text mining, ontology, data integration, machine learning, and information
    architecture. A superior “quantitative ability,” as Gentleman calls it,
    which covers a range of statistical capabilities, is a must, as are
    overarching computing skills. These include core programming abilities,
    such as coding in C++ or Java, or scripting in PERL or Python, says Van
    Criekinge. It is vital to be able to navigate operating systems like UNIX
    and Linux as well as have knowledge of common tools such as Hadoop and NoSQL
     databases, adds Mohan. Experience in data visualization and building
     effective user interfaces, as well as familiarity with hardware, buttresses
      your marketability.
\end{block}
\end{frame}


\begin{frame}
\titlepage
\end{frame}
\begin{frame}[t]\frametitle{Outline}
\tableofcontents[hideallsubsections]
\end{frame}

\section{Perl Modules}

\subsection{Usage of Perl Modules}

\begin{frame}
  \frametitle{Web Scraping in Perl}
  \begin{block}{LWP}
    The libwww-perl collection is a set of Perl modules which provides a simple
    and consistent application programming interface (API) to the World-Wide Web.

    https://metacpan.org/pod/LWP
  \end{block}
\end{frame}

\begin{frame}
  \frametitle{Example: lwp.pl}
\begin{block}{Files}
  \begin{itemize}
    \item pubmedids.txt: A list of pubmed ids
    \item lwp.pl: Get publication titles of these ids from PubMed.
  \end{itemize}
\end{block}
\end{frame}

\subsection{Management of Perl Modules}

\begin{frame}
  \frametitle{Searching}
  \centerline{www.CPAN.org}
\end{frame}

\begin{frame}
  \frametitle{Installation}
\begin{columns}
  \begin{column}{.5\textwidth}
    \begin{itemize}
      \item cpan command: Linux, Mac OS and Windows(Strawberry)
      \item perl -MCPAN -e shell
      \item cpanm
      \item from source
    \end{itemize}
  \end{column}
  \begin{column}{.5\textwidth}
    \begin{itemize}
      \item ppm: ActivePerl
    \end{itemize}
  \end{column}
\end{columns}
\end{frame}



\subsection{Development of Perl Modules}

\begin{frame}
  \frametitle{Development}
  see PUBMED.pm and run.pl
\end{frame}

\subsection{Submit your modules to CPAN}

\begin{frame}
  \frametitle{Why?}
  \begin{itemize}
    \item Easy the installation of your module
    \item Share your work with the community
    \item Increase your importance in the community
    \item Get response form the community to improve your module
  \end{itemize}
\end{frame}

\begin{frame}
  \frametitle{How?}
  http://www.cpan.org/modules/04pause.html
\end{frame}

\section{SNP Annotation}

\begin{frame}
  \frametitle{SNP Annotation}
  \begin{block}{Task}
    Given a list of genome positions, add corrsponding gene symbol to each
    position.
  \end{block}
  \begin{block}{Files}
    \begin{itemize}
      \item Example input file: pos.txt
      \item Annotation database: refGene.txt
      \item Example script: pos_annotate.pl
      \item Optimized Script: pos_annotateV2.pl
    \end{itemize}
  \end{block}
\end{frame}

\subsection{A Perl Script for SNP Annotation}
\begin{frame}
  \frametitle{pos_annotate.pl}
\end{frame}

\begin{frame}
  \frametitle{pos_annotateV2.pl}
\end{frame}

\subsection{SNPAnno module}

\section{Bioinformatics Project in Perl}
\subsection{Cirocs}

\section{BioPerl}

\end{document}
